% ======================================================================
\chapter{Git e GitHub CLI: configuração e login}
\label{cap:git}
% ======================================================================

\section{O que é o Git}

\begin{conceito}[Git]
O Git é um \textbf{sistema distribuído de controle de versão}, gratuito e de
código aberto. Ele registra cada alteração num \textbf{commit}: uma fotografia
do projeto com autor, data e uma mensagem explicando o porquê.
\enquote{Distribuído} quer dizer que cada pessoa tem o histórico completo no
próprio computador.
\end{conceito}

Um arquivo passa por três lugares:

\begin{table}[H]
  \centering
  \caption{As três áreas do Git}
  \label{tab:tres_areas}
  \begin{tabularx}{\textwidth}{@{}l X l@{}}
    \toprule
    \textbf{Área} & \textbf{O que é} & \textbf{Como chega lá} \\
    \midrule
    Diretório de trabalho & os arquivos como você os vê e edita & editando \\
    Área de preparação (\emph{staging}) & o que vai entrar no próximo commit & \comando{git add} \\
    Repositório & o histórico de commits, na pasta oculta \arq{.git} & \comando{git commit} \\
    \bottomrule
  \end{tabularx}
\end{table}

\section{O que é o GitHub CLI}

O GitHub CLI (\texttt{gh}) faz pelo terminal o que você faria no site: criar
repositórios, abrir e revisar Pull Requests, acompanhar workflows. Ele também
resolve a parte mais chata para iniciantes: o \textbf{login}.

% ----------------------------------------------------------------------
\section{Dizendo ao Git quem é você}
% ----------------------------------------------------------------------

Todo commit leva nome e e-mail. Configure uma vez por computador
(Figura~\ref{fig:git_base_config}):

\begin{bashcode}
git config --global user.name "Seu Nome"
git config --global user.email "seu.email@exemplo.com"
\end{bashcode}

\begin{figure}[H]
  \centering
  \includegraphics[width=0.9\textwidth]{./assets/images/14_git_base_config.png}
  \caption{Configurando nome e e-mail}
  \label{fig:git_base_config}
\end{figure}

\begin{dica}[E-mail cadastrado e privacidade]
Use um e-mail cadastrado na sua conta do GitHub, para os commits aparecerem
ligados ao seu perfil. Se não quiser expor o e-mail pessoal, marque
\menu{Keep my email addresses private} em \menu{Settings} \(\rightarrow\)
\menu{Emails} e use o endereço \texttt{...@users.noreply.github.com} que o
GitHub mostra nessa mesma tela.
\end{dica}

Três ajustes que evitam surpresas:

\begin{bashcode}
# novos repositórios começam na branch main, como no GitHub
git config --global init.defaultBranch main

# git pull faz merge quando as branches divergem
git config --global pull.rebase false

# (opcional) VS Code como editor das mensagens de commit
git config --global core.editor "code --wait"
\end{bashcode}

\begin{aviso}[Por que o pull.rebase false]
Sem essa configuração, um \comando{git pull} com alterações dos dois lados
termina em \texttt{fatal: Need to specify how to reconcile divergent
branches}. O instalador do Git para Windows já define essa opção; no Linux e
no macOS é preciso configurar.
\end{aviso}

Para conferir tudo: \comando{git config --global --list}.

% ----------------------------------------------------------------------
\section{Fazendo login no GitHub}
% ----------------------------------------------------------------------

Desde agosto de 2021, o GitHub não aceita a senha da conta em operações do Git
pela linha de comando. É preciso um token ou uma chave SSH. O jeito mais
simples de configurar isso é o GitHub CLI.

\subsection{Pelo GitHub CLI (recomendado)}

\begin{bashcode}
gh auth login
\end{bashcode}

Responda às perguntas assim
(Figuras~\ref{fig:gh_auth_login_0} a~\ref{fig:gh_auth_login_4}):

\begin{table}[H]
  \centering
  \caption{Respostas do \texttt{gh auth login}}
  \label{tab:gh_auth}
  \begin{tabularx}{\textwidth}{@{}X l@{}}
    \toprule
    \textbf{Pergunta} & \textbf{Resposta} \\
    \midrule
    Where do you use GitHub? & \texttt{GitHub.com} \\
    What is your preferred protocol for Git operations? & \texttt{HTTPS} \\
    Authenticate Git with your GitHub credentials? & \texttt{Yes} \\
    How would you like to authenticate GitHub CLI? & \texttt{Login with a web browser} \\
    \bottomrule
  \end{tabularx}
\end{table}

O terminal mostra um código de uso único. Aperte \tecla{Enter} para abrir o
navegador, entre na conta, digite o código e autorize.

\begin{figure}[H]
  \centering
  \includegraphics[width=0.85\textwidth]{./assets/images/13_0_gh_auth_login.png}
  \caption{\texttt{gh auth login}: escolhendo o GitHub.com}
  \label{fig:gh_auth_login_0}
\end{figure}
\begin{figure}[H]
  \centering
  \includegraphics[width=0.85\textwidth]{./assets/images/13_1_gh_auth_login.png}
  \caption{\texttt{gh auth login}: escolhendo o protocolo}
  \label{fig:gh_auth_login_1}
\end{figure}
\begin{figure}[H]
  \centering
  \includegraphics[width=0.85\textwidth]{./assets/images/13_2_gh_auth_login.png}
  \caption{\texttt{gh auth login}: autenticando o Git com as credenciais do GitHub}
  \label{fig:gh_auth_login_2}
\end{figure}
\begin{figure}[H]
  \centering
  \includegraphics[width=0.85\textwidth]{./assets/images/13_3_gh_auth_login.png}
  \caption{\texttt{gh auth login}: login pelo navegador}
  \label{fig:gh_auth_login_3}
\end{figure}
\begin{figure}[H]
  \centering
  \includegraphics[width=0.85\textwidth]{./assets/images/13_4_gh_auth_login.png}
  \caption{\texttt{gh auth login}: código de uso único}
  \label{fig:gh_auth_login_4}
\end{figure}

\begin{conceito}[O que o \enquote{Yes} faz]
Responder \texttt{Yes} em \enquote{Authenticate Git with your GitHub
credentials?} faz o \texttt{gh} emprestar o login dele ao Git. É isso que faz
o \comando{git push} funcionar sem pedir senha. Se você respondeu
\texttt{No}, rode \comando{gh auth setup-git} para configurar depois
\cite{gh_manual}.
\end{conceito}

Para conferir: \comando{gh auth status}.

\subsection{Por chave SSH (alternativa)}

Com SSH, o computador prova quem é usando um par de chaves: a
\textbf{privada} fica só com você, e a \textbf{pública} é cadastrada no GitHub.

\begin{bashcode}
# cria o par de chaves (aceite o local padrão e escolha uma senha)
ssh-keygen -t ed25519 -C "seu.email@exemplo.com"

# cadastra a chave pública na sua conta
gh ssh-key add ~/.ssh/id_ed25519.pub --title "Meu notebook"

# testa a conexão
ssh -T git@github.com
\end{bashcode}

Na primeira conexão, o SSH pergunta se confia no servidor: responda
\texttt{yes}. A resposta esperada é:

\begin{saida}
Hi seu-usuario! You've successfully authenticated, but GitHub does not provide shell access.
\end{saida}

Com SSH, os endereços dos repositórios passam a ser
\texttt{git@github.com:usuario/repositorio.git}, em vez de
\texttt{https://github.com/...}.

\subsection{Por token de acesso pessoal (para scripts)}

Um \textbf{Personal Access Token} (PAT) é uma senha de uso restrito, com prazo
de validade, que você gera no site. Para uso pessoal no terminal, o
\texttt{gh auth login} já resolve. O PAT faz sentido para scripts e
automações.

\begin{enumerate}
  \item Acesse \url{https://github.com/settings/personal-access-tokens/new}
        (token \textbf{fine-grained}).
  \item Dê um nome, uma validade e escolha \textbf{só os repositórios} e
        \textbf{só as permissões} necessárias.
  \item Gere e copie o valor. Ele não aparece de novo.
\end{enumerate}

\begin{aviso}[Token é senha]
Não compartilhe o token nem o grave em arquivos do projeto. Evite o
\texttt{git config credential.helper store}: ele salva o token em texto puro
no arquivo \arq{\textasciitilde/.git-credentials}. O Git Credential Manager
(Windows) e o \texttt{gh auth setup-git} guardam as credenciais de forma
segura.
\end{aviso}

% ----------------------------------------------------------------------
\section{Comandos do dia a dia}
% ----------------------------------------------------------------------

\begin{table}[H]
  \centering
  \caption{Comandos básicos do Git}
  \label{tab:git_basico}
  \begin{tabularx}{\textwidth}{@{}l X@{}}
    \toprule
    \textbf{Comando} & \textbf{O que faz} \\
    \midrule
    \comando{git init} & transforma a pasta atual num repositório \\
    \comando{git clone <url>} & baixa uma cópia completa de um repositório \\
    \comando{git status} & mostra o que mudou e o que está preparado \\
    \comando{git add <arquivo>} & prepara um arquivo para o próximo commit \\
    \comando{git commit -m "..."} & registra as alterações preparadas \\
    \comando{git log --oneline} & mostra o histórico, um commit por linha \\
    \comando{git switch <branch>} & troca de branch \\
    \comando{git switch -c <branch>} & cria uma branch e já entra nela \\
    \comando{git restore <arquivo>} & descarta alterações não preparadas \\
    \comando{git pull} & baixa e integra as novidades do remoto \\
    \comando{git push} & envia os seus commits para o remoto \\
    \comando{git merge <branch>} & traz as alterações de outra branch \\
    \bottomrule
  \end{tabularx}
\end{table}

\begin{table}[H]
  \centering
  \caption{Comandos básicos do GitHub CLI}
  \label{tab:gh_basico}
  \begin{tabularx}{\textwidth}{@{}l X@{}}
    \toprule
    \textbf{Comando} & \textbf{O que faz} \\
    \midrule
    \comando{gh auth login} & faz o login \\
    \comando{gh auth status} & mostra com qual conta você está logado \\
    \comando{gh repo create} & cria um repositório no GitHub \\
    \comando{gh repo clone <dono>/<repo>} & clona um repositório pelo nome \\
    \comando{gh repo fork <dono>/<repo>} & cria um fork na sua conta \\
    \comando{gh pr create} & abre um Pull Request da branch atual \\
    \comando{gh pr list} & lista os Pull Requests abertos \\
    \comando{gh run list} & mostra as execuções do GitHub Actions \\
    \bottomrule
  \end{tabularx}
\end{table}

As listas completas estão nos Apêndices~\ref{ap:git}
e~\ref{ap:gh}.

\begin{dica}[switch e checkout]
Tutoriais antigos usam \comando{git checkout} para trocar e criar branches.
Ele ainda funciona, mas mistura funções demais. O \comando{git switch} (para
branches) e o \comando{git restore} (para arquivos) foram criados para
separar essas tarefas e são mais seguros para quem está começando.
\end{dica}
